\section{Broader applications and research hypotheses}
\label{sec:limitations}

This appendix retains the broader research program. Its hypotheses are not consequences of the reconstruction theorem. The formal framework is exact, deterministic, static, and intentionally bounded. The limitations below define both the assumptions required for broader applications and the research program that follows.

\subsection{Set-based empirical equivalence and exact equality}

Empirical equivalence is equality of reachable observation sets. Real measurements include probabilities, frequencies, noise, finite precision, and time dependence. Two models can have the same support and different probability laws. A stochastic extension should assign measures or kernels to latent and observation spaces. Approximate equivalence should replace exact fibers by metric neighborhoods, couplings, or reconstruction-error bounds.

\subsection{Static interfaces and unrestricted hidden factors}

A dynamical trajectory can distinguish states that one instantaneous output cannot. Future formalization must include state evolution, inputs, intervention, and complete traces. The hidden-factor construction is also deliberately broad. Domain physics, causal structure, conservation laws, geometry, regularity, or minimality may exclude many extensions. The theorem states only that reachable-set equality does not exclude them by itself.

\subsection{Composite chains and local loss analysis}

The conceptual architecture contains apparatus, sampling, encoding, biological access, internal representation, report, and feedback. The current formal core composes these stages into one map. A richer theory should represent each stage and determine where a distinction is preserved, transformed, compressed, dispersed, or collapsed. Such a theory would permit local refinement of one stage and would show whether a later result contains new information or only a re-encoding of an earlier accessible state.

\subsection{The embedded self-interface}

The present analyst stands outside the declared map. Wiener, von Foerster, Maturana and Varela, Rosen, von Neumann, and recent computational observer programs supply distinct precedents for treating observation, self-production, and feedback as coupled rather than external operations \parencite{Wiener1948,VonFoerster1981,MaturanaVarela1980,Rosen1991,VonNeumann1966,Wolfram2023ObserverTheory,RicklesEtAl2023}. The next major formal problem is an embedded system \(S\) that contains a self-model and uses the output of a self-interface to update its own dynamics:
\[
S_t\xrightarrow{\pi_{\mathrm{self}}}O^{\mathrm{self}}_t
\xrightarrow{u}S_{t+1}.
\]
The question is not merely whether \(\pi_{\mathrm{self}}\) is injective. It is whether the system can justify the interface and update rule using only resources generated inside the same process. Coalgebra, domain theory, second-order cybernetics, and active inference are candidate tools \parencite{VonFoerster1981,Abramsky2025,Friston2025ScaleFree}. No general self-justification theorem is claimed here.

\subsection{Cross-domain candidate invariant}

The candidate invariant should be tested separately in logic, computation, cognition, language, biology, artificial intelligence, and physics:
\begin{quote}
\emph{No sufficiently expressive representational system can completely justify or perceive the boundary that generates its own representational capacity.}
\end{quote}
A successful program must define \emph{system}, \emph{representation}, \emph{boundary}, \emph{justify}, and \emph{complete} in each domain, then identify a common theorem or explain why the analogy fails. G\"odel, Tarski, Lawvere, cybernetics, generalized-observer proposals, and self-referential biology supply test cases rather than proof of universality \parencite{Godel1931,Tarski1936,Lawvere1969,Elshatlawy2025,Abramsky2025}.

\subsection{Paradox as a boundary signal}

\begin{hypothesis}[Boundary-signal hypothesis]
A recurrent paradox may sometimes mark a representational system encountering the limit of its own closure rather than an ordinary contradiction inside the model.
\end{hypothesis}
The hypothesis requires a classifier that distinguishes inconsistent assumptions, semantic ambiguity, formal undecidability, empirical underdetermination, and changed access conditions. It is weakened or rejected if all candidate cases are explained by ordinary error without residual structure. It gains support only if a formally defined boundary measure predicts where otherwise independent paradox-like failures occur.

\subsection{Observer-relative mathematics}

A further philosophical hypothesis is that mathematics arises through interaction between an observer and reality, while shared mathematics records structures that remain invariant across different representational routes. In its strongest informal form, ``what we call mathematics is what observers can preserve and agree upon.'' This is not implied by the formal theorem. Mathematical structuralism and category-theoretic invariance provide relevant comparison points \parencite{Shapiro1997,Leinster2014}. Hilbertian axiomatics separates a structure from one concrete realization, while Newtonian and Einsteinian physics illustrate how stable law or invariant relation can survive changes of coordinates and representation \parencite{Hilbert1902,Newton1999,Einstein1905}. Feynman's account of physical law supplies a further methodological reminder that mathematical form earns physical force only through contact with experiment \parencite{Feynman1965}. A future paper would need to distinguish conventional notation, structural equivalence, proof-theoretic invariance, and observer-dependent access to mathematical objects.

\subsection{Materialism, constraint, and unified-state formation}

\begin{hypothesis}[Constraint and unified-state hypothesis]
A physical account of consciousness is incomplete unless it explains how high-dimensional activity is constrained into a stable, integrated, reportable regime rather than merely listing the participating components.
\end{hypothesis}
This hypothesis is not the claim that matter cannot produce consciousness. It asks which dynamical constraints, integration mechanisms, or model-selection processes make a system behave as one temporally coherent state. Hard determinist arguments emphasize the causal closure of biological history, while compatibilist accounts locate agency at the level of reasons, alternatives, control, or second-order endorsement \parencite{Sapolsky2023,List2019,Dennett2003,Frankfurt1971}. These positions do not settle consciousness, but they sharpen the requirement that constraint, agency, and reportability be defined at the correct level. Current active-inference, global-workspace, and integrated-information programs offer distinct answers and are now being subjected to prospective comparison \parencite{SethBayne2022,Cogitate2025,Laukkonen2025,Whyte2026,Robinson2025}. The interface analysis adds an identifiability question: do the proposed measurements distinguish the relevant internal mechanisms, or do multiple mechanisms remain in one empirical class?

A disconfirming outcome would be a sufficiently complete component-level account that predicts unified reportability without requiring an additional system-level constraint variable. A supporting result would require a defined constraint measure that predicts transitions in reportability across sleep, anesthesia, perturbation, or pathology better than competing explanations.

\subsection{Scale-conditioned death and observability}

\begin{hypothesis}[Scale-transition hypothesis of death]
Death is not annihilation of informational structure but a transition in which that structure is redistributed across scales and degrees of freedom, making it no longer observable as a coherent system at the original scale.
\end{hypothesis}
This is a proposed research stance, not a theorem. It must distinguish at least four possibilities: physical persistence of matter and energy, persistence of specific correlations, persistence of organism-level organization, and persistence of subject-level identity. Conservation of energy does not by itself imply conservation of organismal organization or recoverable personal information.

A valid study would define the original-scale coherence criterion, measure selected correlations before and after organismal death, vary spatial and temporal resolution, and determine whether apparent loss at one scale becomes structured signal at another. Failure to detect a structure cannot establish annihilation; detection of a transformed signal cannot establish continuity of the original organism without a reconstruction criterion.

\subsection{Light as an informational standard hypothesis}

\begin{hypothesis}[Light and living organization]
``Light is the informational standard of the living; when it turns into something else it is the informational standard of the dead.''
\end{hypothesis}
The statement is retained exactly as a named hypothesis because it motivated a proposed measurement route. It is not used as a physical fact. To become testable, \emph{informational standard} must be defined as a measurable statistic or family of statistics, and \emph{turns into something else} must specify energy transfer, spectral change, temporal correlation, spatial dispersion, or coupling to another degree of freedom.

Ultraweak photon emission provides one concrete but limited observable. Living and dead mice show measurable differences in ultraweak emission, and emissions in DNA, cultured cells, and stressed plants can be measured under controlled conditions \parencite{PietruszkaMarzec2024,Salari2025,Mould2024}. These observations neither establish a universal living/dead standard nor show that photon emission encodes complete organismal information. A future study should test whether emission statistics predict defined organization or viability measures beyond temperature, oxidative state, metabolic activity, and tissue composition.

\subsection{Single-observer, intentional-forgetting, and external-source hypotheses}

The broader metaphysical hypothesis can be stated without converting it into a conclusion: one foundational observer generates a contained reality, relinquishes control or memory, and experiences apparently independent beings through the system. Intentional forgetting is proposed as a condition that permits uncertainty, love, suffering, discovery, and nontrivial choice. A related external-source hypothesis asks whether structure outside the accessible system could leave internal traces through boundary conditions, invariants, or failures of closure.

These hypotheses are neither proved nor refuted by the present theorem. Their rule of admissibility is strict:
\[
\boxed{\text{No external-source hypothesis earns support without a discriminating prediction.}}
\]
A changed subjective sense of access, revelation, or certainty is not by itself a discriminating prediction. A valid test must specify an outcome unlikely under the physical, plural-observer, information-only, and cognitive alternatives. If no such outcome can be stated, the hypothesis remains metaphysical and must not be used to explain the formal result.

\subsection{Containment, internal reality, and meaning}

The interpretive application is that containment does not imply falseness. Experiences and relations may be real within a containing process even when their ontology is dependent. Meaning, beauty, freedom, and truth may be internally stable outputs of a system without being established as absolute outside it. This claim is philosophical rather than empirical. Leibnizian perspectivism, Whiteheadian relation, and Peircean fallibilism provide three different historical comparisons for this claim \parencite{Leibniz1989,Whitehead1978,Peirce1877,Peirce1878}. Future work should formalize it as invariance under a declared family of transformations: which relations remain stable across observers, interfaces, encodings, or scales?

\subsection{Identity, factorial zero, layered tensors, and geometry}

A separate mathematical note is reserved for the recurring motif that the unit value \(1\) appears as \(0!\), the empty product, \(e^0\), \(\Gamma(1)\), \(\det I\), probability normalization, and identity action. The proposed conceptual sequence is
\[
\text{preserved unit}\longrightarrow
\text{metric structure}\longrightarrow
\text{connection}\longrightarrow
\text{curvature}.
\]
The intended intuition is that Christoffel symbols arise when an identity-like metric structure varies across coordinates, and curvature arises from derivatives and products of the connection. For a specified smooth metric, the Levi-Civita connection depends on the metric and its first derivatives; curvature depends on the connection and its derivatives. Coordinate-dependent metric components can produce nonzero connection coefficients even in flat space, so they do not alone establish curvature. This geometric calculation must be developed independently \parencite{Nakahara2003}. It is not a result of the present theorem, and \(0!=1\) does not by itself generate a connection or curvature tensor.

\subsection{Neural manifolds and reportability}

Published reconstructions of grid-cell population activity motivate an application of the multi-stage interface problem \parencite{Gardner2022,Hermansen2024,DiSarra2025}. The source state, electrode signal, spike-sorted representation, behavioral alignment, dimensional reconstruction, topological summary, biological access, and verbal report are not one object. Future work should define the maps between these levels and test which topological and dynamical conclusions survive justified changes in sampling, preprocessing, reconstruction, and behavioral context.


\subsection{Heterogeneous computational and engineered interfaces}

The same abstract result should not be assumed to have the same meaning in a digital computer, analog device, quantum processor, brain--computer interface, biological observer, multisensor suite, guidance system, or particle detector. Each system has a different physical state space, encoding, admissible operation, error model, and readout. Future formal work should therefore represent a family of interface implementations
\[
(S_i,\eta_i,U_i,\rho_i)
\]
and ask when two implementations realize the same abstract map, when they are refinements of one another, and when an output equivalence conceals different internal processes.

Safety-critical systems make the issue concrete. A multisensor suite combines partial, noisy, and sometimes conflicting measurements into an estimated state or decision; the fusion output remains conditional on coverage, calibration, timing, uncertainty, and the selected fusion architecture \parencite{HallLlinas1997,DurrantWhyteHenderson2008}. An inertial guidance system acts from estimated position, velocity, attitude, and uncertainty rather than direct access to an external ground truth \parencite{Kalman1963,TittertonWeston2004}. A particle-physics exclusion contour is conditional on detector acceptance, reconstruction, the tested model, and a statistical rule \parencite{Cowan2011,ParticleDataGroup2026}. These cases should be formalized as different interface classes rather than treated as intuitive analogies alone.

\subsection{Technical development program}

The next mathematical developments are:
\begin{enumerate}[label=(\roman*)]
\item stochastic observation channels and equality of induced measures;
\item approximate equivalence and quantitative reconstruction error;
\item dynamical observability with inputs and interventions;
\item composition of multiple interfaces and stage-specific fiber analysis;
\item information measures for quotient size, dispersion, and refinement cost;
\item criteria for minimal or canonical representatives;
\item embedded self-interfaces with feedback and model revision;
\item automated countermodel and discriminating-experiment construction;
\item formal comparison with Leibnizian indiscernibility and perspectival representation;
\item computational-irreducibility and bounded-observer criteria;
\item artificial-system interfaces and machine-consciousness indicator tests;
\item translation maps among digital, analog, quantum, neural, and human computational interfaces;
\item state-estimation systems whose actions depend on an internal model and quantified error;
\item multisensor suites that combine partial observations into a bounded state estimate or decision;
\item particle-physics acceptance, blind-region, and exclusion analyses as distinct kinds of boundary claim; and
\item domain-specific applications with independently justified maps.
\end{enumerate}

\subsection{Hypothesis and application register}

\begingroup
\small
\begin{longtable}{@{}>{\raggedright\arraybackslash}p{0.18\textwidth}>{\raggedright\arraybackslash}p{0.17\textwidth}>{\raggedright\arraybackslash}p{0.26\textwidth}>{\raggedright\arraybackslash}p{0.29\textwidth}@{}}
\caption{Status and next requirement for each broader application.}\label{tab:hypothesisregister}\\
\toprule
Program & Current status & Required next object & What would count as progress \\
\midrule
\endfirsthead
\toprule
Program & Current status & Required next object & What would count as progress \\
\midrule
\endhead
Embedded boundary invariant & Cross-domain conjecture & Domain-specific definitions and a common translation theorem & A theorem recurring across at least two nontrivially different domains, or a counterexample that narrows the claim \\
Paradox as boundary signal & Hypothesis & Classifier separating error, undecidability, and empirical underdetermination & Prospective prediction of boundary failures not explained by ordinary inconsistency \\
Observer-relative mathematics & Philosophical hypothesis & Invariance class across formal languages or observers & A precise theorem connecting translation-invariant structure to interface equivalence \\
Constraint and unified consciousness & Empirical and computational program & Defined constraint or integration variable and competing models & Preregistered predictions across state transitions that separate mechanisms \\
Scale-conditioned death & Empirical hypothesis & Coherence criterion across scales and time & Measured transformed structure with a justified continuity or reconstruction relation \\
Light as informational standard & Not yet operationalized & Spectral, temporal, spatial, or information statistic & Predictive value beyond metabolic, thermal, and oxidative confounds \\
Single foundational observer & Metaphysical hypothesis & Unique discriminating prediction & An outcome not shared by plural-observer, physical, information-only, or cognitive alternatives \\
External source or revelation & Metaphysical hypothesis & Publicly checkable boundary condition or invariant & Reproducible evidence with lower likelihood under internal alternatives \\
Identity-to-geometry motif & Separate mathematical note & Independent differential-geometric derivation & Correct theorem showing the precise relation, or a clear statement of the analogy's limits \\
Leibnizian perspective and indiscernibility & Historical-formal comparison & Translation between all properties and interface-accessible predicates & A theorem specifying when interface equivalence preserves or fails Leibnizian indiscernibility \\
Process and constraint dynamics & Dynamical application program & State evolution, constraints, and measurable order parameters & A model separating process, constraint, and observation effects with falsifiable predictions \\
Computational irreducibility and observer theory & Comparative formal program & Explicit coarse-graining and resource bounds & A theorem relating computational aggregation to interface refinement or a counterexample separating them \\
Artificial-system observer interfaces & Empirical and ethical application & Declared outputs, hidden states, self-monitoring channels, and indicator criteria & Tests that distinguish report equivalence from mechanism or consciousness claims \\
Neural manifold and reportability & Empirical application & Explicit maps from neural activity through measurement and reconstruction to report & Results stable across justified preprocessing, sampling, reconstruction, and behavioral changes \\
Heterogeneous physical computation & Formal and engineering application & Encoding, evolution, and readout maps for each implementation & A theorem stating when digital, analog, quantum, neural, or human interfaces realize the same abstract relation \\
Inertial navigation and guidance & State-estimation application & Dynamical model, sensor map, error state, and uncertainty relation & Correct predictions of when distinct physical trajectories remain indistinguishable to the estimator \\
Multisensor suite and data fusion & Safety-critical state-estimation application & Sensor-specific observation maps, timing, uncertainty, and a declared fusion rule & Correctly calibrated estimates whose blind regions, residuals, and unresolved state components are explicit \\
Particle-physics forbidden and excluded regions & Experimental application & Separate definitions for theoretical prohibition, detector inaccessibility, and statistical exclusion & Reproducible limits whose dependence on model, acceptance, and confidence rule is explicit \\
\bottomrule
\end{longtable}
\endgroup

The distinction among theorem, computation, interpretation, empirical application, and metaphysical hypothesis must remain explicit.
